\documentclass[a4paper]{article}

\usepackage{graphicx, amssymb, amsmath, textcomp, booktabs}
\usepackage[libertine,vvarbb]{newtxmath}
\usepackage[scr=rsfso]{mathalfa}
\usepackage[T1]{fontenc}
\usepackage{minted}
\usepackage{listings, color, setspace, titlesec, fancyhdr, mdframed}
\usepackage{ucharclasses}
\usepackage{xunicode, xltxtra}
\usepackage[inner=1.25cm, outer=0.8cm, top=1.7cm, bottom=1.5cm]{geometry}
\usepackage{tocloft}
\usepackage{nameref}
\usepackage{verbatim}
\usepackage{relsize}
\usepackage{fontspec}
\usepackage[colorlinks, linkcolor = black]{hyperref}
\usepackage[table]{xcolor}
\usepackage{tabularx}
\usepackage{enumitem}
\usepackage{multicol}
\setlength{\columnsep}{16pt}

\usepackage{xeCJK}
\setCJKmainfont{Source Han Serif SC}[Scale=0.8, BoldFont=Source Han Serif SC Bold]
\setCJKmonofont{SimHei}[Scale=0.8]
\setCJKsansfont{KaiTi}[Scale=0.8]

% 为封面定义黑体（纯 xeCJK 方式）
\setCJKfamilyfont{zhhei}{SimHei}

\setmainfont{Linux Libertine O}[Scale=0.925]
\setmonofont{Consolas}[Scale=0.775]

\XeTeXlinebreaklocale "zh"
\XeTeXlinebreakskip = 0pt plus 1pt

\setlength{\parindent}{0em}\setlength{\parskip}{1pt}
\setlength\itemsep{1pt}

% header
\pagestyle{fancy}
\setlength{\headsep}{0.1cm}
\chead{Aeyone}
\rhead{Page \thepage}
\lhead{Template}

% toc
\renewcommand\cftsecfont{\Large}

% section style
\titleformat{\section}
{\huge\bfseries}
{\thesection.}
{4pt}
{}
\titleformat{\subsection}
{\large\bfseries}
{\thesubsection.}
{4pt}
{}
\titlespacing{\section}{0pt}{0pt}{0pt}
\titlespacing{\subsection}{0pt}{0pt}{0pt}
\titlespacing{\subsubsection}{0pt}{0pt}{0pt}

% minted line number style
\renewcommand{\theFancyVerbLine}{\sffamily \textcolor[rgb]{0.5,0.5,0.5}{\scriptsize {\arabic{FancyVerbLine}}}}

% minted 全局默认（所有语言生效）
\setminted{
    style=xcode,
    linenos,
    baselinestretch=0.8,
    tabsize=3,
    fontsize=\normalsize,
    frame=single,
    framesep=1mm,
    framerule=0.3pt,
    numbersep=1mm,
    breaklines=true,
    breaksymbolsepleft=2pt,
    breakbytoken=false,
	showtabs=true,
    tab={\relscale{0.6} $\big\vert \ \ \ $ \relscale{1}},
}

\begin{document}

% === 封面 ===
\hypersetup{pageanchor=false}
\begin{titlepage}
\thispagestyle{empty}
\vspace*{\fill}
\begin{center}
    {\CJKfamily{zhhei}\Huge 算法竞赛模板}\\[1cm]
    {\Large Aeyone}\\[0.5cm]
    {\large \today}
\end{center}
\vspace*{\fill}
\end{titlepage}
\hypersetup{pageanchor=true}

\tableofcontents
\newpage
\begin{multicols}{2}
\begin{spacing}{0.6}

\section{一些杂碎}

\subsection{主模板}
\inputminted{cpp}{\detokenize{Template/一些杂碎/01 - main.txt}}

\subsection{run.sh}
\begin{minted}{bash}
#!/bin/bash
touch $1.in $1.out
g++ -std=c++23 -Wall -Wextra $1.cpp -o out.o
./out.o < $1.in > $1.out
\end{minted}

\subsection{对拍器}
\inputminted{cpp}{\detokenize{src/Debug/run.cpp}}

\columnbreak

\subsection{随机化}
\begin{minted}[breakanywhere=true]{cpp}
static mt19937_64 rng(chrono::steady_clock::now().time_since_epoch().count());
#define rand(l, r) (uniform_int_distribution<long long>((l),(r))(rng))
shuffle(a.begin(), a.end(), rng) // 打乱一个数组
\end{minted}

\subsection{\texorpdfstring{线性求 \(1 \sim n\) 逆元}{线性求 1 至 n 逆元}}
\inputminted{cpp}{\detokenize{Template/一些杂碎/02B - 线性求1~n逆元.txt}}

\subsection{进制转换}
\inputminted{cpp}{\detokenize{Template/一些杂碎/03 - 进制转换.txt}}

\subsection{int128 输出流自定义}
\inputminted{cpp}{\detokenize{Template/一些杂碎/04 - int128 输出流自定义.txt}}

\end{spacing}
\end{multicols}

\begin{spacing}{0.6}
\newpage
\section{图论}

\subsection{强连通分量缩点（SCC）}
\inputminted{cpp}{\detokenize{Template/图论/01 - 强连通分量缩点（SCC）.txt}}

\newpage
\subsection{树链剖分（HLD）}
\inputminted{cpp}{\detokenize{Template/图论/02 - 树链剖分（HLD）.txt}}

\newpage
\subsection{最大流（Dinic）}
\inputminted{cpp}{\detokenize{Template/图论/03 - 最大流（Dinic）.txt}}

\newpage
\subsection{费用流（最小费用流）}
\inputminted{cpp}{\detokenize{Template/图论/04 - 费用流（最小费用流）.txt}}

\newpage
\section{字符串}

\subsection{KMP}
\inputminted{cpp}{\detokenize{Template/字符串/01 - KMP.txt}}

\subsection{马拉车（Manacher）}
\inputminted{cpp}{\detokenize{Template/字符串/02 - 马拉车（Manacher）.txt}}

\subsection{Trie}
\inputminted{cpp}{\detokenize{Template/字符串/03 - Trie.txt}}

\subsection{字符串哈希}
\inputminted{cpp}{\detokenize{Template/字符串/04 - 字符串哈希.txt}}

\newpage

\end{spacing}

\begin{multicols}{2}
\begin{spacing}{0.6}
\section{数学}

\subsection{快速幂}
\inputminted{cpp}{\detokenize{Template/数学/01A - 快速幂.txt}}

\subsection{矩阵快速幂}
\inputminted{cpp}{\detokenize{Template/数学/01B - 矩阵快速幂.txt}}
\columnbreak

\subsection{组合计数}
\inputminted{cpp}{\detokenize{Template/数学/02 - 组合计数.txt}}
\columnbreak

\subsection{根号求欧拉函数}
\inputminted{cpp}{\detokenize{Template/数学/03A - 根号求欧拉函数.txt}}

\subsection{线性筛求欧拉函数（Euler 筛法）}
\inputminted{cpp}{\detokenize{Template/数学/03B - 线性筛求欧拉函数（Euler筛法）.txt}}

\subsection{拓展欧几里得}
\inputminted{cpp}{\detokenize{Template/数学/04 - 拓展欧几里得.txt}}

\columnbreak

\subsection{高斯消元}
\inputminted{cpp}{\detokenize{Template/数学/05 - 高斯消元.txt}}

\subsection{线性基（Basis）}
\inputminted{cpp}{\detokenize{Template/数学/06 - 线性基（Basis）.txt}}

\end{spacing}
\end{multicols}


\begin{spacing}{0.6}

\newpage
\subsection{计算几何}
\inputminted{cpp}{\detokenize{Template/数学/07 - 计算几何.txt}}

\newpage
\subsection{多项式乘法（NTT）}
\inputminted{cpp}{\detokenize{Template/数学/08- 多项式乘法（NTT）.txt}}

\subsection{米勒-罗宾素数检验（Miller–Rabin）}
\inputminted{cpp}{\detokenize{Template/数学/09 - 米勒-罗宾素数检验（Miller–Rabin）.txt}}
\newpage

\end{spacing}

\begin{multicols}{2}
\begin{spacing}{0.6}

\section{数据结构}

\subsection{并查集（DSU）}
\inputminted{cpp}{\detokenize{Template/数据结构/01 - 并查集（DSU）.txt}}

\subsection{普通莫队}
\inputminted{cpp}{\detokenize{Template/数据结构/02 - 普通莫队.txt}}

\subsection{树状数组（FenwickTree）}
\inputminted{cpp}{\detokenize{Template/数据结构/03 - 树状数组（FenwickTree）.txt}}
\columnbreak

\subsection{ST表（SparseTable）}
\inputminted{cpp}{\detokenize{Template/数据结构/04 - ST表（SparseTable）.txt}}
\columnbreak

\subsection{线段树（区间合并，不带修）}
\inputminted{cpp}{\detokenize{Template/数据结构/05 - 线段树（区间合并，不带修）.txt}}
\columnbreak

\subsection{懒标记线段树（基础区间操作）}
\inputminted{cpp}{\detokenize{Template/数据结构/06A - 懒标记线段树（基础区间操作）.txt}}
\columnbreak

\subsection{懒标记线段树（双重区间操作）}
\inputminted{cpp}{\detokenize{Template/数据结构/06B - 懒标记线段树（双重区间操作）.txt}}
\columnbreak

\subsection{懒标记线段树（动态开点+区间操作）}
\inputminted{cpp}{\detokenize{Template/数据结构/06C - 懒标记线段树（动态开点+区间操作）.txt}}
\columnbreak

\subsection{可持久化线段树（单点操作）}
\inputminted{cpp}{\detokenize{Template/数据结构/07A - 可持久化线段树（单点操作）.txt}}
\columnbreak

\subsection{可持久化线段树（区间操作）}
\inputminted{cpp}{\detokenize{Template/数据结构/07A - 可持久化线段树（区间操作）.txt}}
\columnbreak

\subsection{吉如一线段树（区间求和）}
\inputminted{cpp}{\detokenize{Template/数据结构/08A - 吉如一线段树（区间求和）.txt}}
\columnbreak

\subsection{吉如一线段树（区间求和+区间历史最值）}
\inputminted{cpp}{\detokenize{Template/数据结构/08B - 吉如一线段树（区间求和+区间历史最值）.txt}}

\end{spacing}
\end{multicols}
\end{document}
